\chapter{Conclusiones}

Existen múltiples simuladores diseñados para replicar el comportamiento de los autómatas celulares; sin embargo, fue necesario un simulador específico para este problema. El sistema está formado por dos partes principales: el simulador, que a su vez se divide en autómatas celulares y autómatas celulares neuronales, y el módulo de entrenamiento.

El simulador de autómatas celulares permite evolucionarlos en espacios de $2000 \times 2000$ gracias a la paralelización personalizada de procesos mediante CUDA, su conexión mediante \textit{bindings} de C++ y su manejo de \textit{buffer} para la optimización. Además, este módulo permite definir el autómata celular mediante una vecindad semitotalística (que también permite definir reglas arbitrarias y totalísticas), los colores del simulador, el vector de regla, el control de la evolución y la manipulación directa del espacio.

Por otro lado

Adicionalmente, el módulo de entrenamiento, que fue creciendo conforme a las necesidades del sistema, permite entrenar redes neuronales multicapa de forma totalmente personalizada, mediante comandos o mediante la interfaz gráfica del simulador. Este módulo permite usar múltiples evoluciones como datos y definir todas las capas, las neuronas y la función de activación de cada capa, la codificación de los datos de entrada, el tamaño de la vecindad, el tipo de frontera de los datos y los hiperparámetros del entrenamiento, como las épocas, el tamaño del lote, la pérdida objetivo, la tasa de aprendizaje, la reducción de la tasa de aprendizaje o la semilla aleatoria.

Entrenar y evaluar el desempeño de los autómatas celulares neuronales bajo diferentes configuraciones

Con el simulador fue posible realizar los seis experimentos. Tres de ellos son comportamientos ampliamente conocidos, en los que es posible saber exactamente cuál será el siguiente estado: (1) el juego de la vida, (2) el \textit{sandpile} y (3) la hormiga de Langton. Con ellos pudimos definir configuraciones de los autómatas celulares neuronales y, más específicamente, de la red neuronal, que facilitan o permiten inferir un comportamiento.

Con el juego de la vida se demostró la capacidad inicial del sistema para simular sistemas complejos. El juego de la vida requiere una de las redes neuronales más pequeñas por la poca información que necesita: con tan solo 9 neuronas en la capa de entrada (una por vecino), 1 capa oculta de 4 neuronas con ReLU y 2 neuronas con función Linear en la capa de salida es posible aprender su comportamiento. Este fue el primer hito: permitió seguir adelante y dio paso a definiciones más extensas. Además, este experimento muestra la relación entre las reglas arbitrarias y las semitotalísticas, pues un entrenamiento incompleto, variando la pérdida objetivo, permite que una regla aprenda otro comportamiento.

Otro fenómeno interesante es la estabilidad de la regla B3/S23: al variar la tasa de aprendizaje, únicamente se aprendieron 2 comportamientos: la regla 0, en la que cualquier patrón se vuelve 0, o la misma regla. Lo contrario ocurre con otras reglas, como B3/S238, B38/S238 o B3/S236, que generan vida adicional y que, al variar su pérdida objetivo, sí llegan a aprender otras reglas. Esto puede deberse a que los pesos de la red neuronal se inicializan en cero; sin embargo, cabe la posibilidad de que esta regla sea más estable y, a su vez, una forma común del resto de las evoluciones.

Finalmente, este primer experimento muestra la importancia de la vecindad al entrenar: con una vecindad más pequeña no se puede aprender el comportamiento, pero una vecindad más grande no impide replicarlo.

El juego de la vida es un sistema sencillo porque, al ser semitotalístico, sus condiciones están dadas por valores lineales, por lo que no requiere técnicas adicionales. Lo contrario ocurre con la pila de arena de Bak–Tang–Wiesenfeld, que tiene 17 estados. Este experimento tuvo dos aportaciones principales. La primera fue reducir la tasa de aprendizaje con \textit{plateau} para evitar oscilaciones en la pérdida y llegar a la pérdida objetivo de manera uniforme. La segunda, y la que más mejoró la replicación, fue usar codificación \textit{one-hot} para los estados que no son lineales. Esto no ocurría en B3/S23 porque únicamente tiene 2 estados; sin embargo, incluso con dos estados, codificarlos en \textit{one-hot} tuvo mejor rendimiento que usar una sola entrada por vecino.


El último de estos tres experimentos sirvió para demostrar la capacidad de manejar múltiples estados y escasez de información. La hormiga de Langton es una máquina de Turing en dos dimensiones, por lo que, en un espacio de $n \times m$, únicamente cambian 2 células en cada paso y el resto se mantiene constante.

Con esa información, sin requerir preprocesamiento adicional, el autómata celular neuronal demostró su capacidad de extraer información de comportamientos poco representados, con una alta efectividad incluso usando 10 estados, y fue capaz de persistir por más de 12,000 pasos sin equivocarse.

Queda demostrado que el autómata celular neuronal puede replicar algunos comportamientos complejos, especialmente de configuraciones bien conocidas, pero con restricciones. La mayor limitación es la representación de escenarios: ante patrones que no aparecen en los datos de entrenamiento, la red se equivoca, aunque su comportamiento sigue siendo similar.

Esta demostración dio pie a los siguientes tres experimentos, que pusieron a prueba la capacidad de replicación en entornos reales: (1) la propagación de fuego con vista nadir (FLAME~3), (2) el movimiento de tormentas (SEVIR) y (3) el crecimiento de moho en una caja de Petri.

Resulta complicado, o casi imposible, determinar la función que rige cada uno de estos sistemas. Por esto se usan los autómatas celulares: para determinar cómo las relaciones locales definen el comportamiento colectivo. En estos casos, cómo el fuego se propaga y quema diferentes zonas según la geografía, cómo las tormentas se forman, desaparecen y se mueven, o cómo el moho se propaga por un espacio.

Sin embargo, cada experimento presentó sus limitaciones.

La propagación del fuego se puede replicar, pero existen múltiples elementos que no se consideran y que limitan la capacidad de replicación; por ejemplo, se ignoran el relieve, la temperatura, el clima, las carreteras, los árboles y el viento. Ignorar cada una de estas condiciones por separado no sería un problema, pero al ignorarlas todas, la abstracción dentro del simulador pierde tanto contexto sobre el sistema que la replicación no es muy fiel. El autómata celular neuronal fue capaz de replicar el comportamiento de propagación; sin embargo, la velocidad y la dirección dependen exclusivamente del conjunto de datos con el que fue entrenado.

Además, este sistema se vio limitado por la cantidad de información del conjunto de datos. Al tener únicamente dos grabaciones pequeñas, los comportamientos aprendidos dependen mucho de las condiciones en las que se grabaron: en la primera evolución, el sistema se calienta de manera uniforme y tiene un salto conforme se acerca la barrera de fuego; en la segunda, el sistema está frío y las temperaturas solo aumentan cuando llega la barrera de fuego. Son dos comportamientos diferentes.

Por otro lado, la simulación del movimiento de tormentas presenta otro inconveniente. El sistema es capaz de replicar el movimiento, pero únicamente en las secuencias que se mueven en una sola dirección (no en las que se mueven en remolino); en las demás, desplaza las tormentas y elimina sus otros comportamientos. El único comportamiento que se aprendió fue el desplazamiento de las tormentas en una sola dirección, hacia el este (característico del viento en Estados Unidos). Una posible mejora sería cambiar el autómata celular determinista por uno probabilístico: el ruido podría ayudar a definir ciertos comportamientos y aplicarlo no es computacionalmente complejo (ni modificarlo en el simulador), pero sí es un cambio en la definición de autómata celular propuesta, lo que sería una línea de estudio completamente diferente. Por último, lo que hace que este sistema no sea apto para simularse con un autómata celular neuronal es que no es un sistema cerrado. Como se mencionó en los capítulos previos, un autómata celular es un sistema cerrado: a partir de una configuración inicial, evoluciona sin que factores externos intervengan en su comportamiento. De lo contrario, la red asumiría cada factor externo como propio del sistema. Por esta razón se propone el último experimento, que consiste en grabaciones del moho \textit{Physarum polycephalum} en una caja de Petri, un sistema cerrado que tiene una configuración inicial.



